\chapter{Some results in comparison geometry}

\section*{B.1. The Gauss Lemma and radial vector fields}

\begin{lemma}[Gauss Lemma]
\label{appb-lem-gauss-lemma}
\source{chow-knopf-ricci-flow:page-0300}
\dcref{Lemma B.1}
\group{local-comparison-geometry}

Let $(\mathcal M^n,g)$ be a complete Riemannian manifold, let $p\in\mathcal M^n$, and let $V,W\in T_p\mathcal M^n$ satisfy $\langle V,W\rangle=0$. Then along the geodesic ray determined by $V$, the differential of the exponential map preserves orthogonality to the radial direction:
\[
\langle (\exp_p)_*(V_{tV}),(\exp_p)_*(W_{tV})\rangle = 0.
\]
\end{lemma}

\begin{proof}
\source{chow-knopf-ricci-flow:page-0300}
\sketch This is the standard Gauss lemma for the exponential map.
\end{proof}

\begin{corollary}[Gradient of the radial distance]
\label{appb-cor-gradient-radial-distance}
\source{chow-knopf-ricci-flow:page-0301, chow-knopf-ricci-flow:page-0302}
\uses{appb-lem-gauss-lemma}
\dcref{Corollary B.2}
\group{local-comparison-geometry}


For sufficiently small $\epsilon>0$, on $\exp_p(B(\vec 0,\epsilon))\setminus\{p\}$ the gradient of the distance function from $p$ agrees with the radial unit vector field:
\[
\nabla f = \frac{\partial}{\partial r}.
\]
\end{corollary}

\begin{proof}
\source{chow-knopf-ricci-flow:page-0301, chow-knopf-ricci-flow:page-0302}
\sketch Decompose an arbitrary tangent vector into radial and orthogonal parts and apply the Gauss lemma to show that the distance function has derivative $1$ in the radial direction and vanishing derivative in orthogonal directions.
\end{proof}

\section*{B.2. Conjugate points}

\begin{definition}[Conjugate point]
\label{appb-def-conjugate-point}
\source{chow-knopf-ricci-flow:page-0302}
\dcref{Definition B.3}
\group{local-comparison-geometry}

A point $q\in\mathcal M^n$ is a conjugate point of $p\in\mathcal M^n$ if $q$ is a critical value of the exponential map $\exp_p:T_p\mathcal M^n\to\mathcal M^n$.
\end{definition}

\begin{definition}[Conjugate radius]
\label{appb-def-conjugate-radius}
\source{chow-knopf-ricci-flow:page-0303}
\uses{appb-def-conjugate-point}
\dcref{Definition B.4}
\group{local-comparison-geometry}


The conjugate radius of a point $p\in\mathcal M^n$ is
\[
r_c(p)\coloneqq \sup\{r:(\exp_p)_*\text{ is nonsingular on }B(\vec 0,r)\}.
\]
\end{definition}

\begin{remark}[Immersion at the conjugate radius]
\label{appb-rem-conjugate-radius-immersion}
\source{chow-knopf-ricci-flow:page-0303}
\uses{appb-def-conjugate-radius}
\dcref{Remark B.5}
\group{local-comparison-geometry}


The restriction $\exp_p|_{B(\vec 0,r_c(p))}$ is an immersion.
\end{remark}

\begin{lemma}[Conjugate points in limits of geodesic $2$-gons]
\label{appb-lem-limit-geodesic-2gon-conjugate}
\source{chow-knopf-ricci-flow:page-0303, chow-knopf-ricci-flow:page-0304}
\uses{appb-def-conjugate-point}
\dcref{Lemma B.6}
\group{local-comparison-geometry}


Let $(\mathcal M^n,g)$ be complete and let $\{\Gamma_i\}$ be a sequence of nondegenerate proper geodesic $2$-gons with vertices and initial tangent data converging to limits $p_\infty$, $q_\infty$, $V_\infty$, and $\ell_\infty$. Then $\ell_\infty>0$, and $q_\infty$ is a conjugate point to $p_\infty$ along the limit geodesic determined by $V_\infty$.
\end{lemma}

\begin{proof}
\source{chow-knopf-ricci-flow:page-0303, chow-knopf-ricci-flow:page-0304}
\sketch If $q_\infty$ were not conjugate to $p_\infty$, then the inverse function theorem would give a neighborhood on which the exponential map is injective for all nearby basepoints and tangent vectors. For sufficiently large $i$, this would force the two distinct geodesic segments of $\Gamma_i$ joining the same endpoints to have the same initial vector, a contradiction.
\end{proof}

\begin{remark}[Failure in varying manifolds]
\label{appb-rem-limit-failure-varying-manifolds}
\source{chow-knopf-ricci-flow:page-0304}
\dcref{Remark B.7}
\group{local-comparison-geometry}

The conclusion of Lemma~\ref{appb-lem-limit-geodesic-2gon-conjugate} can fail for sequences of varying manifolds, even with uniformly bounded geometry.
\end{remark}

\begin{example}[Collapsing flat tori]
\label{appb-ex-collapsing-flat-tori}
\source{chow-knopf-ricci-flow:page-0304}
\dcref{Example B.8}
\group{local-comparison-geometry}

For collapsing flat tori, a sequence of nondegenerate proper geodesic $2$-gons can converge to a geodesic segment with no conjugate points.
\end{example}

\begin{example}[Spherical orbifold example]
\label{appb-ex-spherical-orbifold}
\source{chow-knopf-ricci-flow:page-0304}
\dcref{Example B.9}
\group{local-comparison-geometry}

On the quotients of the round sphere by cyclic rotation, two distinct minimizing geodesics can converge to a single geodesic segment with no conjugate points.
\end{example}

\section*{B.3. Jacobi fields and focal points}

\begin{definition}[Jacobi field]
\label{appb-def-jacobi-field}
\lean{ChowKnopf.IsGeodesicVariationOn, ChowKnopf.IsJacobiFieldOn}
\leanok
\source{chow-knopf-ricci-flow:page-0305}
\dcref{Definition B.10}
\group{local-comparison-geometry}

A Jacobi field is a variation vector field of a one-parameter family of geodesics.
\end{definition}

\begin{lemma}[Conjugate points and the differential of $\exp_p$]
\label{appb-lem-conjugate-differential-exp}
\source{chow-knopf-ricci-flow:page-0305}
\uses{appb-def-conjugate-point}
\dcref{Lemma B.11}
\group{local-comparison-geometry}


A point $\exp_p(V)$ is conjugate to $p$ if and only if there exists a nonzero $W\in T_p\mathcal M^n$ orthogonal to $V$ such that $(\exp_p)_*(W_V)=0$.
\end{lemma}

\begin{proposition}[Jacobi fields characterize conjugate points]
\label{appb-prop-jacobi-characterizes-conjugacy}
\source{chow-knopf-ricci-flow:page-0305}
\uses{appb-def-conjugate-point, appb-def-jacobi-field}
\dcref{Proposition B.12}
\group{local-comparison-geometry}


Let $(\mathcal M^n,g)$ be complete and connected. A point $q$ is conjugate to $p$ along a geodesic $\gamma$ if and only if there exists a nonzero Jacobi field $J$ along $\gamma$ with $J(0)=J(1)=0$.
\end{proposition}

\begin{proof}
\source{chow-knopf-ricci-flow:page-0305}
\sketch The Jacobi field obtained by varying the initial velocity of a geodesic gives the forward direction, and the converse is immediate from the characterization of critical points of the exponential map.
\end{proof}

\begin{lemma}[The Jacobi equation]
\label{appb-lem-jacobi-equation}
\source{chow-knopf-ricci-flow:page-0305, chow-knopf-ricci-flow:page-0306}
\uses{appb-def-jacobi-field}
\dcref{Lemma B.13}
\group{local-comparison-geometry}


A field $J$ along a geodesic with unit tangent vector $T$ is a Jacobi field if and only if
\[
\nabla_T\nabla_T J = R(T,J)T.
\]
\end{lemma}

\begin{corollary}[Orthogonal decomposition of Jacobi fields]
\label{appb-cor-jacobi-orthogonal-decomposition}
\source{chow-knopf-ricci-flow:page-0305, chow-knopf-ricci-flow:page-0306}
\uses{appb-lem-jacobi-equation}
\dcref{Corollary B.14}
\group{local-comparison-geometry}


Every Jacobi field $J$ along a geodesic with unit tangent vector $T$ decomposes uniquely as
\[
J = J_0 + (at+b)T,
\]
where $a,b\in\mathbb R$ and $\langle J_0,T\rangle\equiv 0$.
\end{corollary}

\begin{proof}
\source{chow-knopf-ricci-flow:page-0306}
\sketch The component of $J$ along the geodesic direction has vanishing second derivative, so it must be affine in the arclength parameter.
\end{proof}

\begin{definition}[Normal bundle]
\label{appb-def-normal-bundle}
\source{chow-knopf-ricci-flow:page-0306}
\dcref{Definition B.15}
\group{local-comparison-geometry}

If $l:\Sigma^k\hookrightarrow \mathcal M^n$ is an immersion, the normal bundle $N\Sigma^k$ is the subbundle of $l^*T\mathcal M^n$ whose fiber at $p\in\Sigma^k$ consists of vectors orthogonal to $l_*(T_p\Sigma^k)$.
\end{definition}

\begin{definition}[Focal point]
\label{appb-def-focal-point}
\source{chow-knopf-ricci-flow:page-0306}
\uses{appb-def-normal-bundle}
\dcref{Definition B.18}
\group{local-comparison-geometry}


If $l:\Sigma^k\hookrightarrow\mathcal M^n$ is an immersion, a point $q\in\mathcal M^n$ is a focal point of $l(\Sigma^k)$ if it is a critical value of the normal exponential map $\exp_{\mid N\Sigma^k}$.
\end{definition}

\section*{B.4. The Rauch comparison theorem}

\begin{theorem}[Rauch comparison theorem]
\label{appb-thm-rauch-comparison}
\source{chow-knopf-ricci-flow:page-0307}
\uses{appb-def-jacobi-field}
\dcref{Theorem B.20}
\group{comparison-theorems}


Let $(\mathcal M^n,g)$ and $(\overline{\mathcal M}^n,\bar g)$ be complete Riemannian manifolds and let $\gamma,\bar\gamma$ be unit-speed geodesics of the same length. If every sectional curvature along $\gamma$ is bounded above by the corresponding sectional curvature along $\bar\gamma$, then the norm of any Jacobi field along $\gamma$ dominates that of the corresponding Jacobi field along $\bar\gamma$, provided the initial data match and $\bar\gamma$ has no conjugate point on the interval.
\end{theorem}

\begin{corollary}[Comparison under nonpositive model curvature]
\label{appb-cor-rauch-nonpositive-model}
\source{chow-knopf-ricci-flow:page-0307}
\uses{appb-thm-rauch-comparison}
\dcref{Corollary B.21}
\group{comparison-theorems}


If $(\mathcal M^n,g)$ is complete and $\sec(g)\le K$ for some $K>0$, then a Jacobi field along a geodesic satisfies the standard upper comparison estimate with the model space of curvature $K$.
\end{corollary}

\begin{corollary}[Injectivity radius bound under an upper sectional curvature bound]
\label{appb-cor-rauch-injectivity-radius}
\source{chow-knopf-ricci-flow:page-0307}
\uses{appb-thm-rauch-comparison}
\dcref{Corollary B.22}
\group{comparison-theorems}


Under an upper sectional curvature bound by $K>0$, the exponential map is an immersion on the open ball of radius $\pi/\sqrt K$ in each tangent space.
\end{corollary}

\begin{proof}
\source{chow-knopf-ricci-flow:page-0307}
\sketch This is the standard consequence of Rauch comparison for the model sphere of constant curvature $K$.
\end{proof}

\section*{B.5. The Laplacian comparison theorem}

\begin{theorem}[Laplacian comparison]
\label{appb-thm-laplacian-comparison}
\source{chow-knopf-ricci-flow:page-0307, chow-knopf-ricci-flow:page-0313, chow-knopf-ricci-flow:page-0314}
\uses{appb-cor-gradient-radial-distance}
\dcref{Theorem B.23}
\group{comparison-theorems}


If the Ricci curvature is bounded below by $(n-1)K$, then the Laplacian of the distance function from a point is bounded above by the corresponding radial Laplacian in the simply connected space form of curvature $K$.
\end{theorem}

\begin{corollary}[Model-space radial comparison]
\label{appb-cor-model-space-radial-comparison}
\source{chow-knopf-ricci-flow:page-0313, chow-knopf-ricci-flow:page-0314}
\uses{appb-thm-laplacian-comparison}
\dcref{Corollary B.24}
\group{comparison-theorems}


The radial comparison inequality specializes to the familiar Euclidean, spherical, and hyperbolic model formulas according to the sign of $K$.
\end{corollary}

\section*{B.6. The Toponogov comparison theorem}

\begin{theorem}[Toponogov comparison theorem]
\label{appb-thm-toponogov-comparison}
\source{chow-knopf-ricci-flow:page-0315, chow-knopf-ricci-flow:page-0316}
\uses{appb-thm-rauch-comparison, appb-thm-laplacian-comparison}
\dcref{Theorem B.40}
\group{comparison-theorems}


Let $(\mathcal M^n,g)$ be complete with sectional curvature bounded below by $H\in\mathbb R$. Then geodesic triangles in $\mathcal M^n$ are no thinner than the corresponding comparison triangles in the simply connected space form of curvature $H$.
\end{theorem}

\begin{corollary}[Angle comparison]
\label{appb-cor-angle-comparison}
\source{chow-knopf-ricci-flow:page-0316}
\uses{appb-thm-toponogov-comparison}
\dcref{Corollary B.41}
\group{comparison-theorems}


Under the hypotheses of Toponogov comparison, corresponding angles in geodesic triangles satisfy the usual lower comparison bound.
\end{corollary}

\begin{lemma}[Total convexity of half spaces]
\label{appb-lem-halfspace-totally-convex}
\source{chow-knopf-ricci-flow:page-0316, chow-knopf-ricci-flow:page-0317}
\uses{appb-thm-toponogov-comparison}
\dcref{Lemma B.42}
\group{comparison-theorems}


For any ray $\gamma$, the associated half space $\mathcal H_\gamma$ is closed and totally convex.
\end{lemma}

\begin{lemma}[Exhaustion by half spaces]
\label{appb-lem-halfspace-exhaustion}
\source{chow-knopf-ricci-flow:page-0317, chow-knopf-ricci-flow:page-0318}
\uses{appb-lem-halfspace-totally-convex}
\dcref{Lemma B.43}
\group{comparison-theorems}


The union of the half spaces associated to the shifted rays of a fixed ray exhausts the manifold.
\end{lemma}

\begin{lemma}[Basic Busemann inequalities]
\label{appb-lem-busemann-basic}
\source{chow-knopf-ricci-flow:page-0318, chow-knopf-ricci-flow:page-0319}
\uses{appb-lem-halfspace-exhaustion}
\dcref{Lemma B.44}
\group{comparison-theorems}


For the Busemann function $b_\gamma$ of a ray $\gamma$, one has the standard inequalities comparing $b_\gamma$ to distance and the corresponding shifted-ray descriptions of half spaces.
\end{lemma}

\begin{lemma}[Compactness of geodesic segments in a compact set]
\label{appb-lem-compactness-geodesic-segments}
\source{chow-knopf-ricci-flow:page-0319, chow-knopf-ricci-flow:page-0320}
\dcref{Lemma B.45}
\group{comparison-theorems}

In a complete manifold, the family of geodesic segments with endpoints constrained to a compact set admits convergent subsequences after passing to a subsequence of initial data.
\end{lemma}

\begin{proposition}[Nonconvexity would force a topological contradiction]
\label{appb-prop-nonconvexity-contradiction}
\source{chow-knopf-ricci-flow:page-0320, chow-knopf-ricci-flow:page-0321, chow-knopf-ricci-flow:page-0322}
\uses{appb-thm-toponogov-comparison, appb-lem-halfspace-totally-convex}
\dcref{Proposition B.54}
\group{comparison-theorems}


If $\sec(g)\ge 0$, then the relevant half spaces and Busemann sublevel sets are totally convex, and the proof by contradiction using a minimal geodesic loop closes.
\end{proposition}

\begin{definition}[Shifted ray]
\label{appb-def-shifted-ray}
\lean{ChowKnopf.Ray, ChowKnopf.shiftedRay, ChowKnopf.shiftedRay_apply}
\leanok
\source{chow-knopf-ricci-flow:page-0322}
\dcref{Definition B.55}
\group{comparison-theorems}

If $\gamma$ is a ray and $s\ge 0$, its shifted ray $\gamma_s$ is given by $\gamma_s(s')=\gamma(s+s')$.
\end{definition}

\begin{lemma}[Monotonicity of half spaces]
\label{appb-lem-shifted-ray-monotone}
\source{chow-knopf-ricci-flow:page-0322}
\uses{appb-def-shifted-ray}
\dcref{Lemma B.56}
\group{comparison-theorems}


If $s\le t$, then $\mathbb B_{\gamma_t}\subseteq \mathbb B_{\gamma_s}$.
\end{lemma}

\begin{lemma}[Distance characterization of shifted half spaces]
\label{appb-lem-shifted-ray-distance}
\source{chow-knopf-ricci-flow:page-0322, chow-knopf-ricci-flow:page-0323}
\uses{appb-def-shifted-ray, appb-lem-shifted-ray-monotone}
\dcref{Lemma B.57}
\group{comparison-theorems}


If $0<s<t$, then
\[
\mathbb B_{\gamma_s}=\{x\in\mathcal M^n:d(x,\mathbb B_{\gamma_t})<t-s\}.
\]
\end{lemma}

\begin{definition}[Busemann sublevel set]
\label{appb-def-busemann-sublevel}
\lean{ChowKnopf.raysFrom, ChowKnopf.busemannHalfSpace,
  ChowKnopf.busemannSublevelSet, ChowKnopf.mem_busemannSublevelSet_iff}
\leanok
\source{chow-knopf-ricci-flow:page-0323}
\dcref{Definition B.58}
\group{comparison-theorems}

Given an origin $O\in\mathcal M^n$ and $s\ge 0$, the Busemann sublevel set is
\[
C_s\coloneqq \bigcap_{\gamma\in\mathcal R(O)} \mathbb B_{\gamma_s}.
\]
\end{definition}

\begin{lemma}[Busemann sublevel sets and the Busemann function]
\label{appb-lem-busemann-sublevels}
\source{chow-knopf-ricci-flow:page-0323, chow-knopf-ricci-flow:page-0324}
\uses{appb-def-busemann-sublevel}
\dcref{Lemma B.59}
\group{comparison-theorems}


For the Busemann function $\bar b$ associated to $O$, one has $C_s=\{x:\bar b(x)\le s\}$. In particular, the sets $C_s$ are nested and their boundaries are the level sets of $\bar b$.
\end{lemma}

\begin{corollary}[Closed ball containment]
\label{appb-cor-busemann-ball-containment}
\source{chow-knopf-ricci-flow:page-0324}
\uses{appb-lem-busemann-sublevels}
\dcref{Corollary B.60}
\group{comparison-theorems}


If $s>0$, then $C_s$ contains the closed metric ball of radius $s$ centered at $O$.
\end{corollary}

\begin{corollary}[Exhaustion by Busemann sublevel sets]
\label{appb-cor-busemann-exhaustion}
\source{chow-knopf-ricci-flow:page-0324}
\uses{appb-lem-busemann-sublevels}
\dcref{Corollary B.61}
\group{comparison-theorems}


For any origin $O$, the union $\bigcup_{s\ge 0} C_s$ is all of $\mathcal M^n$.
\end{corollary}

\begin{proposition}[Compact and totally convex Busemann sublevel sets]
\label{appb-prop-busemann-sublevel-compact-convex}
\source{chow-knopf-ricci-flow:page-0324}
\uses{appb-lem-busemann-sublevels, appb-thm-toponogov-comparison}
\dcref{Proposition B.62}
\group{comparison-theorems}


For any origin $O$ and any $s\ge 0$, the set $C_s$ is compact and totally convex.
\end{proposition}

\begin{corollary}[Lower bound for the Busemann function]
\label{appb-cor-busemann-below}
\source{chow-knopf-ricci-flow:page-0324}
\uses{appb-prop-busemann-sublevel-compact-convex}
\dcref{Corollary B.63}
\group{comparison-theorems}


The Busemann function associated to an origin is bounded below.
\end{corollary}

\begin{proposition}[Parallel level sets]
\label{appb-prop-busemann-level-sets-parallel}
\source{chow-knopf-ricci-flow:page-0325}
\uses{appb-prop-busemann-sublevel-compact-convex}
\dcref{Proposition B.64}
\group{comparison-theorems}


For $s<t$, the sublevel set $C_s$ is exactly the set of points in $C_t$ at distance at least $t-s$ from $\partial C_t$.
\end{proposition}

\section*{B.7. Estimating injectivity radius in positive curvature}

\begin{theorem}[Injectivity radius lower bound]
\label{appb-thm-inj-radius-positive-curvature}
\source{chow-knopf-ricci-flow:page-0325, chow-knopf-ricci-flow:page-0326, chow-knopf-ricci-flow:page-0327, chow-knopf-ricci-flow:page-0328}
\uses{appb-prop-busemann-sublevel-compact-convex, appb-prop-busemann-level-sets-parallel}
\dcref{Theorem B.65}
\group{comparison-theorems}


If $(\mathcal M^n,g)$ is complete, noncompact, has positive sectional curvature, and satisfies $\sec(g)\le K$ for some $K>0$, then
\[
\operatorname{inj}(\mathcal M^n,g)\ge \frac{\pi}{\sqrt K}.
\]
\end{theorem}

\begin{proof}
\source{chow-knopf-ricci-flow:page-0327, chow-knopf-ricci-flow:page-0328}
\sketch Assuming the injectivity radius bound fails, one produces a short smooth geodesic loop in a compact totally convex Busemann sublevel set and then contradicts minimality by a second variation argument along a minimizing geodesic from the loop to a point on a ray.
\end{proof}

\begin{definition}[Proper geodesic $k$-gon]
\label{appb-def-proper-geodesic-k-gon}
\source{chow-knopf-ricci-flow:page-0325, chow-knopf-ricci-flow:page-0326}
\dcref{Definition B.66}
\group{comparison-theorems}

For $k\in\mathbb N^*$, a proper geodesic $k$-gon is a cyclic collection of unit-speed geodesic segments joining $k$ pairwise distinct vertices. It is nondegenerate if all interior angles are nonzero.
\end{definition}

\begin{definition}[Geodesic $k$-gon]
\label{appb-def-geodesic-k-gon}
\source{chow-knopf-ricci-flow:page-0326}
\uses{appb-def-proper-geodesic-k-gon}
\dcref{Definition B.67}
\group{comparison-theorems}


A geodesic $k$-gon is a proper geodesic $j$-gon for some $j\le k$.
\end{definition}

\begin{lemma}[Compactness of the minimizing polygon space]
\label{appb-lem-polygon-space-compact}
\source{chow-knopf-ricci-flow:page-0326, chow-knopf-ricci-flow:page-0327}
\uses{appb-def-proper-geodesic-k-gon, appb-def-geodesic-k-gon}
\dcref{Lemma B.68}
\group{comparison-theorems}


The space $\Lambda$ of nondegenerate geodesic $2$-gons in a fixed compact Busemann sublevel set with bounded total length is closed in the ambient compact parameter space $\Xi$, hence compact.
\end{lemma}

\begin{lemma}[A minimizing polygon is smooth]
\label{appb-lem-minimizing-polygon-smooth}
\source{chow-knopf-ricci-flow:page-0327}
\uses{appb-lem-polygon-space-compact}
\dcref{Lemma B.69}
\group{comparison-theorems}


If $\beta\in\Lambda$ minimizes total length, then $\beta$ is a smooth geodesic loop.
\end{lemma}

\begin{proof}
\source{chow-knopf-ricci-flow:page-0327, chow-knopf-ricci-flow:page-0328}
\sketch If a minimizing polygon were not smooth at a vertex, a one-parameter shortening deformation would have negative first variation, contradicting minimality.
\end{proof}

\begin{example}[Unique focal point of a Euclidean sphere]
\label{appb-ex-euclidean-sphere-focal-point}
\dcref{Example B.19}
\group{local-comparison-geometry}
\uses{appb-def-focal-point}
\source{chow-knopf-ricci-flow:page-0306}
For the round sphere $S_r^{n-1}(p)\subset\mathbb R^n$, the unique focal point is the center $p$.
\end{example}

\begin{example}[Trivial normal bundle in codimension $0$]
\label{appb-ex-normal-bundle-codim-zero}
\dcref{Example B.16}
\group{local-comparison-geometry}
\uses{appb-def-normal-bundle}
\source{chow-knopf-ricci-flow:page-0306}
If $\Sigma^n=\mathcal M^n$, then $N\Sigma^n$ is naturally identified with $\Sigma^n$ and the normal exponential map is the identity.
\end{example}

\begin{example}[Normal bundle of a point]
\label{appb-ex-normal-bundle-point}
\dcref{Example B.17}
\group{local-comparison-geometry}
\uses{appb-def-normal-bundle}
\source{chow-knopf-ricci-flow:page-0306}
If $\Sigma^0=\{p\}$, then $N\{p\}=T_p\mathcal M^n$ and the normal exponential map is the ordinary exponential map $\exp_p$.
\end{example}
